\section{Detection Engineering Fundamentals}

Detection engineering focuses on identifying the behaviors, patterns, and signals that indicate malicious 
activity within an environment. While threat hunting is exploratory and hypothesis-driven, detection engineering 
is structured and repeatable. The goal is to translate knowledge of adversary behavior into reliable, high-fidelity 
detections that operate continuously in a SIEM such as Splunk or Microsoft Sentinel. This chapter introduces the 
foundational concepts and provides practical SPL and KQL examples that reflect real-world detection workflows.

% --------------------------------------------------------------------
% SECTION 1: Behavioral Detection Philosophy
% --------------------------------------------------------------------

\subsection*{Behavioral Detection Philosophy}

A good detection begins with understanding the behavior you want to identify. Instead of relying on static indicators 
such as IP addresses or file hashes, modern detections focus on techniques—how an adversary executes code, moves laterally, 
steals credentials, or evades defenses. These behaviors tend to be more stable over time and more resilient to simple 
changes in tooling.

For example:

\begin{itemize}
    \item A malware family may change its command-and-control servers weekly.
    \item But its use of encoded PowerShell commands often remains consistent.
    \item A threat actor may rotate infrastructure.
    \item But its parent-child process chains remain predictable.
\end{itemize}

Behavior-based detections are durable, transferable, and harder for adversaries to evade.

% --------------------------------------------------------------------
% SECTION 2: Suspicious Command-Line Activity
% --------------------------------------------------------------------

\subsection*{Suspicious Command-Line Activity}

One of the most common detection patterns involves monitoring for suspicious command-line activity. PowerShell is frequently 
abused by adversaries because it is built into Windows and provides extensive system access. The following example detects 
encoded PowerShell commands, which are often used to hide malicious intent.

\textbf{SPL Example}

\begin{lstlisting}
index=sysmon EventCode=1 Image="*powershell.exe"
| search CommandLine="*EncodedCommand*"
| table _time, Computer, User, ParentImage, Image, CommandLine
\end{lstlisting}

\textbf{KQL Equivalent}

\begin{lstlisting}
Sysmon
| where EventID == 1 and Image endswith "powershell.exe"
| where CommandLine contains "EncodedCommand"
| project TimeGenerated, Computer, User, ParentImage, Image, CommandLine
\end{lstlisting}

This pattern is foundational for detecting obfuscated execution, malware loaders, and post-exploitation frameworks.

% --------------------------------------------------------------------
% SECTION 3: Parent-Child Process Anomalies
% --------------------------------------------------------------------

\subsection*{Parent-Child Process Anomalies}

Legitimate processes tend to follow predictable execution chains, while malicious activity often introduces unusual or rare 
combinations. The following example identifies uncommon parent-child pairs.

\textbf{SPL Example}

\begin{lstlisting}
index=sysmon EventCode=1
| stats count by ParentImage, Image
| where count < 3
\end{lstlisting}

\textbf{KQL Equivalent}

\begin{lstlisting}
Sysmon
| where EventID == 1
| summarize Count=count() by ParentImage, Image
| where Count < 3
\end{lstlisting}

This technique is effective for detecting:

\begin{itemize}
    \item LOLBins used for execution
    \item malware spawning unusual child processes
    \item privilege escalation chains
    \item lateral movement tools
\end{itemize}

% --------------------------------------------------------------------
% SECTION 4: Authentication-Based Detections
% --------------------------------------------------------------------

\subsection*{Authentication-Based Detections}

Authentication anomalies often indicate credential misuse or lateral movement. Excessive failed logons, unusual logon types, 
or logons from unexpected hosts are strong behavioral signals.

\textbf{SPL Example}

\begin{lstlisting}
index=wineventlog EventCode=4625
| stats count by AccountName, WorkstationName
| where count > 10
\end{lstlisting}

\textbf{KQL Equivalent}

\begin{lstlisting}
SecurityEvent
| where EventID == 4625
| summarize Count=count() by Account, Workstation
| where Count > 10
\end{lstlisting}

This pattern supports detections for:

\begin{itemize}
    \item brute-force attempts
    \item password spraying
    \item compromised accounts
    \item lateral movement
\end{itemize}

% --------------------------------------------------------------------
% SECTION 5: Network-Based Detections
% --------------------------------------------------------------------

\subsection*{Network-Based Detections}

Network telemetry helps identify command-and-control activity, data exfiltration, or lateral movement. Rare outbound 
connections are often a strong signal of malicious behavior.

\textbf{SPL Example}

\begin{lstlisting}
index=sysmon EventCode=3
| stats count by DestinationIp, DestinationPort
| where count < 5
\end{lstlisting}

\textbf{KQL Equivalent}

\begin{lstlisting}
Sysmon
| where EventID == 3
| summarize Count=count() by DestinationIp, DestinationPort
| where Count < 5
\end{lstlisting}

This pattern is useful for detecting:

\begin{itemize}
    \item unusual outbound traffic
    \item malware beaconing
    \item lateral movement tunnels
    \item exfiltration channels
\end{itemize}

% --------------------------------------------------------------------
% SECTION 6: Tuning and Validation
% --------------------------------------------------------------------

\subsection*{Tuning and Validation}

Effective detection engineering requires tuning and validation. A detection should be tested against real telemetry to 
ensure it produces meaningful results without overwhelming analysts with noise.

Tuning often involves:

\begin{itemize}
    \item adding filters for known-good activity
    \item whitelisting legitimate processes or hosts
    \item enriching events with additional context
    \item adjusting thresholds based on environment baselines
\end{itemize}

Detections should be reviewed periodically as adversary behavior evolves or as new telemetry sources become available.

% --------------------------------------------------------------------
% SECTION 7: Why Detection Engineering Matters
% --------------------------------------------------------------------

\subsection*{Why Detection Engineering Matters}

The examples in this chapter represent the core building blocks of detection engineering: identifying suspicious behavior, 
correlating related events, and transforming threat intelligence into actionable, automated logic. These fundamentals form 
the basis for building reliable detections, dashboards, and playbooks that reflect real-world adversary techniques. Later 
chapters expand these concepts into full detection patterns and advanced detection engineering workflows.
